\section[几何 Satake]{几何 Satake（1 月 22 日）}

\srcpage{113}

\subsection{关于 $D_{\text{ét}}(\Bun_G)$ 的更多内容}

迄今为止，对于通常的 $G/E$，我们已经定义了 $\Perf_{\ol{\Fq}}$ 上的一个「Artin $v$-叠」$\Bun_G$，以及 étale 层的导出范畴
\[
D_{\text{ét}}(\Bun_G,\mathbb Z/\ell^n\mathbb Z).
\]

\begin{question}[Drinfeld]
能否定义一个范畴 $D(\Bun_G,\mathbb Z[1/p])$，使得
\[
D(\Bun_G,\mathbb Z[1/p])\otimes_{\mathbb Z[1/p]}\mathbb Z/\ell^n\mathbb Z=D_{\text{ét}}(\Bun_G,\mathbb Z/\ell^n\mathbb Z)?
\]
\end{question}

通常人们会通过动机理论来处理这个问题，但动机理论在这里不适用。

由于 $D_{\text{ét}}(\Bun_G,\mathbb Z/\ell^n\mathbb Z)$ 分层成若干片
\[
D_{\text{ét}}(\Bun_G^b,\mathbb Z/\ell^n)\cong D(G_b(E),\mathbb Z/\ell^n),
\]
人们可以期望，假设中的 $D(\Bun_G,\mathbb Z[1/p])$ 也允许这样的分层，分成若干片
\[
D(\Bun_G^b,\mathbb Z[1/p])\cong D(G_b(E),\mathbb Z[1/p]).
\]

部分答案：这样的范畴以 $\mathbb Z_\ell$ 系数存在，进而也以 $\ol{\mathbb Q}_\ell$ 系数存在。但给定 $\iota\colon\ol{\mathbb Q}_\ell\cong\ol{\mathbb Q}_{\ell'}$ 时，是否存在等价
\[
D(\Bun_G,\ol{\mathbb Q}_\ell)\cong D(\Bun_G,\ol{\mathbb Q}_{\ell'}),
\]
尚不清楚。

我们只能典范地用 $\mathbb Z_\ell$ 系数工作，因为该范畴隐含地知道 Tate 扭转
\[
\mathbb Z_\ell(1)\cong\Hom(\mathbb Q_\ell/\mathbb Z_\ell,\ol{\mathbb F}_q^\times),
\]
它是一个秩为 1 的自由 $\mathbb Z_\ell$-模。因此我们需要 $\mathbb Z_\ell(1)$ 上的一个 $\mathbb Z[1/p]$-结构。

例如，我们可以对所有 $\ell\neq p$ 选取同构 $\mathbb Z_\ell(1)\cong\mathbb Z_\ell$。固定这样的选择后，这样的范畴似乎应当存在。一个相关的事实是：在 Langlands 对偶一侧，我们有一个典范的、$\mathbb Z_\ell$ 上的 L-参数 Artin 叠，即
\[
\mathrm{Par}_G(A)=\frac{\{\text{连续 }W_E\to\widehat G(A)\}}{\widehat G-\text{共轭}},\qquad A/\mathbb Z_\ell.
\]
\yizhu{原文左端下标作 $G$，而分子与分母均为 $\widehat G$，疑为 $\mathrm{Par}_{\widehat G}$ 之笔误；图版如此，照录。}{A（原文 p.113）}

这必须在 $\mathbb Z_\ell$ 上定义，因为驯惯性 $\prod_{\ell\neq p}\mathbb Z_\ell(1)$ 只能非平凡地映入 $\mathbb Z_\ell$-代数。然而，固定一个拓扑生成元 $\tau\in\prod_{\ell\neq p}\mathbb Z_\ell(1)$ 后，可以构造 Weil 群的一个部分离散化版本 $W_E^\tau\subset W_E$，把驯惯性替换为 $\mathbb Z[1/p]\cdot\tau$。于是存在一个 $\mathbb Z[1/p]$ 上的 Artin 叠，参数化 $\{W_E^\tau\to\widehat G\}/\widehat G$；它在基变换到 $\mathbb Z_\ell$ 后给出所有典范的此类叠 [DHKM, Zhu20]。

\begin{question}[Drinfeld]
当 $G=\mathrm{SL}_2$ 时，能否把上述导出范畴显式化？
\end{question}

关键问题如下。考虑
\[
\Bun_G^{b_1}\xleftarrow{i_1}\Bun_G\xleftarrow{i_2}\Bun_G^{b_2}.
\]
那么 $i_2^*Ri_{1*}\colon D(G_{b_1}(E),\Lambda)\to D(G_{b_2}(E),\Lambda)$ 是什么？

抽象的回答（来自局部卡片的结构）如下：考虑
\[
\begin{tikzcd}
\widetilde{\mathcal M}_{b_2} \arrow[r] & \Bun_G \\
\widetilde{\mathcal M}_{b_2}^{b_1} \arrow[r,"f"] \arrow[u] & \Bun_G^{b_1} \arrow[u]
\end{tikzcd}
\]
存在 $G_{b_2}(E)$ 在 $\widetilde{\mathcal M}_{b_2}^{b_1}$ 上的作用。对 $\pi\in\Rep(G_{b_2}(E))$，令 $[\pi]$ 为 $\Bun_G^{b_1}$ 上对应的层。那么函子 $i_2^*Ri_{1*}$ 由下式给出（可能差一个平移）
\[
\pi\mapsto R\Gamma(\widetilde{\mathcal M}_{b_2}^{b_1},f^*[\pi]).
\]

\srcpage{114}

\begin{example}
设 $G=\mathrm{SL}_2$，$b_1\leftrightarrow\cO^2$，$b_2\leftrightarrow\cO(-1)\oplus\cO(1)$，则 $\widetilde{\mathcal M}_{b_2}^{b_1}$ 参数化 $\cO(1)$ 被 $\cO(-1)$ 的、且同构于 $\cO^2$ 的扩张。或者，选取一个与 $\cO^2$ 的同构，得到一个覆盖 $\widetilde{\widetilde{\mathcal M}}{}^{b_1}_{b_2}$，它参数化余核为 $\cO(1)$ 的单射 $\cO(-1)\hookrightarrow\cO^2$。
\end{example}

换言之，这些就是饱和单射 $\cO(-1)\hookrightarrow\cO^2$。非饱和的单射可以延拓为映射 $\cO\hookrightarrow\cO^2$。这些容易理解，因为 $\cO$ 的整体截面是 $E$。所以
\[
\widetilde{\widetilde{\mathcal M}}{}^{b_1}_{b_2}=BC(\cO(1))^2\setminus(\mathrm{GL}_2(E)\cdot\Delta BC(\cO(1))).
\]

我们需要计算
\[
R\Gamma(\mathrm{SL}_2(E),\pi\otimes R\Gamma(\widetilde{\widetilde{\mathcal M}}{}^{b_1}_{b_2},\Lambda)).
\]
特别地，我们需要计算 $R\Gamma(\widetilde{\widetilde{\mathcal M}}{}^{b_1}_{b_2},\Lambda)$。由于 $BC(\cO(1))^2\cong\mathbb D^2$，而 $\mathrm{GL}_2(E)\cdot\Delta BC(\cO(1))$ 是 $\mathbb P^1(E)$ 份在 $0$ 处粘合的 $\mathbb D$，可以用切除来分析。这样会得到平凡表示与 Steinberg 表示的若干拷贝，最多差一些扭转。

\begin{example}
设 $b_1\leftrightarrow\cO(-1)\oplus\cO(1)$，$b_2\cong\cO(-2)\oplus\cO(2)$，则 $\widetilde{\widetilde{\mathcal M}}{}^{b_1}_{b_2}$ 参数化饱和单射 $\cO(-2)\hookrightarrow\cO(-1)\oplus\cO(1)$。非饱和的单射可以延拓为 $\cO(-1)\hookrightarrow\cO(-1)\oplus\cO(1)$，所以
\[
\widetilde{\widetilde{\mathcal M}}{}^{b_1}_{b_2}=BC(\cO(1))\times BC(\cO(3))\setminus\bigl(\text{$\ul{E}\times BC(\cO(2))\times BC(\cO(1))$ 的像}\bigr),
\]
其中映射为 $(xz,yz)\leftarrow(x,y,z)$。再次可以用切除来计算。
\end{example}

\subsection{我们要去向何处？}

我们想要提取 L-参数
\[
\varphi\colon W_E\to\widehat G\in H^1(W_E,\widehat G).
\]
我们仍然需要让对偶群 $\widehat G$ 出现。

想法：L-参数是「谱信息」，它作为 Hecke 算子作用在 $D_{\text{ét}}(\Bun_G,\Lambda)$ 上的「特征值」出现。

关键事实是：Hecke 算子由 $\Rep\widehat G$ 枚举。

\begin{warning}
在 $p$ 进群的表示论中，有一个经典的 Hecke 算子概念。那不是我们这里讨论的东西。
\end{warning}

\subsection{Hecke 算子}

\begin{definition}
设 $\mathrm{Hecke}_G$ 是 $\Perf_{\ol{\Fq}}$ 上的小 $v$-叠，其点函子为
\[
\mathrm{Hecke}_G(S)=\left\{(\mathcal E_1,\mathcal E_2,S^\#,f):\ \begin{array}{l}
\mathcal E_1,\mathcal E_2\in\Bun_G(S)\\[2pt]
S^\#\in\Div_X^1(S)\text{ 是 }S\text{ 在 }E\text{ 上的 untilt，模去 Frobenius}\\[2pt]
f\colon\mathcal E_1|_{X_S\setminus S^\#}\cong\mathcal E_2|_{X_S\setminus S^\#}\text{ 在 }S^\#\text{ 处亚纯}
\end{array}\right\}.
\]
\end{definition}

如此定义的 $\mathrm{Hecke}_G$ 是「无限维的」，因为修正 $f$ 不是「有界的」。对 $G$ 的支配余特征 $\mu$，存在子叠 $\mathrm{Hecke}_{G,\le\mu}$ 界住修正的极点，使得映射 $h_1,h_2$ 紧合，并且可由空间菱形表示。
\[
\begin{tikzcd}
& \mathrm{Hecke}_{G,\le\mu} \arrow[dl,"h_1"'] \arrow[dr,"h_2"] & \\
\Bun_G & & \Bun_G\times\Div^1
\end{tikzcd}
\]

\srcpage{115}

于是我们得到形如
\[
Rh_{2*}h_1^*\colon D_{\text{ét}}(\Bun_G,\Lambda)\to D_{\text{ét}}(\Bun_G\times\Div^1,\Lambda)
\]
的算子。由于 $\Div^1\cong(\Spa\breveE)^\circ/W_E$，我们有 $D_{\text{ét}}(\Bun_G\times\Div^1,\Lambda)\cong D_{\text{ét}}(\Bun_G,\Lambda)^{W_E}$。这里我们隐含地使用了 $D_{\text{ét}}(\Bun_G,\Lambda)$ 沿 $(\Spa\breveE)^\circ\to(\Spa\ol{\Fq})^\circ$ 的基变换不变性。

实际上，更好的做法是允许对应上的核。

\begin{theorem}[几何 Satake，第一个版本]
存在一个典范的精确幺半函子
\[
\Rep_{\mathbb Z_\ell}(\widehat G)\to D_{\text{ét}}(\mathrm{Hecke}_G,\Lambda),
\]
记作 $V\mapsto\mathcal S_V$。
\end{theorem}

于是我们得到 Hecke 算子
\[
T_V\colon Rh_{2*}(h_1^*\otimes\mathcal S_V)\colon D_{\text{ét}}(\Bun_G,\Lambda)\to D_{\text{ét}}(\Bun_G,\Lambda)^{W_E}.
\]
映射 $V\mapsto\mathcal S_V$ 的幺半性蕴涵 $T_W\circ T_V\cong T_{V\otimes W}$。

\subsection{经典几何 Satake}

我们从经典设置开始（由 Mirkovic–Vilonen、Lusztig、Ginzburg 等人发展）。

设 $G/\mathbb C$ 是约化群。正环路群 $L^+G$ 的点函子为 $L^+G(A)=G(A[[t]])$。这定义了一个无限维仿射概形。环路群 $LG$ 的点函子为 $LG(A)=G(A((t)))$。它是一个归纳概形。

\begin{definition}
仿射 Grassmann 簇为 $\Gr_G=LG/L^+G$。其点函子把 $A$ 送到 $A[[t]]$ 上、在 $A((t))$ 上平凡化的 $G$-主丛 $\mathcal E$ 的集合。
\end{definition}

这是一个归纳射影概形，过渡映射是闭浸入。

\begin{definition}
定义 $G$ 的 Satake 范畴为 $\mathrm{Sat}_G:=\Perv_{L^+G}(\Gr_G;\mathbb Z)$。
\end{definition}

轨道空间 $L^+G\backslash LG/L^+G$ 可以与 $X_*^+$（$G$ 的支配余特征）等同。这个双射把 $\mu\in X_*^+$ 送到由 $\mu(t)$ 代表的双陪集。$\mu(t)$ 的 $L^+G$-轨道记作 $\Gr_\mu$，其闭包是一个射影概形 $\Gr_{G,\le\mu}\subset\Gr_G$，即一个（仿射）Schubert 簇。

特别地，对每个 $\mu$，我们有相交复形（intersection complex）$IC_\mu=IC_{\Gr_{G,\le\mu}}\in\mathrm{Sat}_G$。按定义，
\[
IC_\mu=\operatorname{Im}\left(
\underbrace{{}^p j_{\mu!}\mathbb Z[d_\mu]}_{\text{「标准」}\Delta_\mu}
\to
\underbrace{{}^p j_{\mu*}\mathbb Z[d_\mu]}_{\text{「余标准」}\nabla_\mu}
\right),
\]
其中 $j_\mu\colon\Gr_\mu\hookrightarrow\Gr_{G,\le\mu}$。这里 $d_\mu:=\dim\Gr_\mu=\langle 2\rho,\mu\rangle$。

\begin{remark}
以 $\mathbb Q$ 为系数时，${}^pj_{\mu!}\mathbb Q[d_\mu]\xrightarrow{\sim}IC_{\mu,\mathbb Q}\xrightarrow{\sim}{}^pj_{\mu*}\mathbb Q[d_\mu]$。
\end{remark}

\begin{remark}
$\Delta_\mu,IC_\mu,\nabla_\mu$ 的结构赋予 $\mathrm{Sat}_G$ 一个「最高权范畴」的结构，其权由 $X_*^+$ 给出。以 $\mathbb Q$ 为系数时，它是半单的。
\end{remark}

\srcpage{116}

\begin{definition}
$\mathrm{Sat}_G$ 上有一个卷积幺半结构，如下给出。考虑交换图
\[
\begin{tikzcd}
(L^+G\backslash LG/L^+G)\times(L^+G\backslash LG/L^+G) \\
L^+G\backslash LG\times^{L^+G}LG/L^+G \arrow[u,"\pi"] \arrow[d,"m"] \\
L^+G\backslash LG/L^+G
\end{tikzcd}
\]
对 $A,B\in\mathrm{Sat}_G$，定义 $A\star B=Rm_*\pi^*(A\boxtimes B)$。
\end{definition}

\begin{theorem}[Mirkovic–Vilonen {[MV07]}]
上同调函子
\[
(\mathrm{Sat}_G,\star)\xrightarrow{\oplus H^i(\Gr_G,-)}(\Vect,\otimes)
\]
是纤维函子，并且 $(\mathrm{Sat}_G,\star)$ 可以提升为对称幺半结构，使得 $\oplus H^i(\Gr_G,-)$ 成为对称幺半函子。对应的 Tannaka 群是 $\widehat G$，所以我们有交换三角形
\[
\begin{tikzcd}
(\mathrm{Sat}_G,\star) \arrow[rr,"\sim"] \arrow[dr,"{\oplus H^i(\Gr_G,-)}"'] & & (\Rep\widehat G,\otimes) \arrow[dl,"\text{遗忘}"] \\
& (\Vect,\otimes) &
\end{tikzcd}
\]
\end{theorem}

\begin{itemize}
\item 以 $\mathbb Q$ 为系数时，$IC_\mu$ 对应 $V_\mu$，即 $\widehat G$ 的权为 $\mu\in X_*^+=X^*_+(\widehat G)$ 的最高权表示。
\item 以 $\mathbb F_p$ 为系数时，$IC_\mu$ 对应不可约表示 $L_\mu$，而 ${}^pj_{\mu!}\mathbb F_p[d_\mu]$ 对应 $\Delta_\mu$，${}^pj_{\mu*}\mathbb F_p[d_\mu]$ 对应 $\nabla_\mu$。
\end{itemize}

\subsection{$p$ 进几何 Satake}

我们想要为 $B^+_{\mathrm{dR}}$-仿射 Grassmann 簇得到这个故事的一个版本。

\begin{definition}
设 $G/E$、$\Div^1$ 如通常。我们定义一个函子 $\Gr_G\to\Div^1$，它参数化 $S^\#\in\Div^1(S)$、$X_S$ 在 $S^\#$ 处的完备化上的一个 $G$-主丛 $\mathcal E$，以及 $(X_S)^\wedge_{S^\#}\setminus S^\#$ 上的一个平凡化。
\end{definition}

\begin{remark}
这里有一个微妙之处：$(X_S)^\wedge_{S^\#}$ 是什么意思？我们只对 $S=\Spa(R,R^+)$ 为 affinoid 的情形定义它。此时 $S^\#=\Spa(R^\#,R^{\#+})$ 也是 affinoid，并且我们有
\[
\theta\colon W_{\cO_E}(R^+)[1/[\varpi]]\twoheadrightarrow R^\#.
\]
我们定义 $B^+_{\mathrm{dR}}(R^\#)$ 为 $W_{\cO_E}(R^+)[1/[\varpi]]$ 的 $(\ker\theta)$-进完备化。（这与任一满足 $\Spa(A,A^+)\supset S^\#$ 的开 affinoid 子集 $\Spa(A,A^+)\subset X_S$ 沿 $S^\#$ 的完备化相同。）可以证明 $\ker\theta=(\xi)$ 是主理想；于是定义 $\mathbb B_{\mathrm{dR}}(R^\#):=B^+_{\mathrm{dR}}(R^\#)[1/\xi]$。我们定义 $(X_S)^\wedge_{S^\#}:=\Spa B^+_{\mathrm{dR}}(R^\#)$，以及 $(X_S)^\wedge_{S^\#}\setminus S^\#:=\Spa\mathbb B_{\mathrm{dR}}(R^\#)$。
\end{remark}

所以 $\Gr_G(S)=\{B^+_{\mathrm{dR}}(R^\#)$ 上的 $G$-主丛，在 $\mathbb B_{\mathrm{dR}}(R^\#)$ 上平凡化$\}$。我们有 $\Gr_G=LG/L^+G$，其中 $LG$ 是把 $R^\#$ 送到 $G(\mathbb B_{\mathrm{dR}}(R^\#))$ 的函子，$L^+G$ 是把 $R^\#$ 送到 $G(B^+_{\mathrm{dR}}(R^\#))$ 的函子。

\begin{definition}
定义局部 Hecke 叠
\[
\mathrm{Hk}_G=L^+G\backslash\Gr_G=L^+G\backslash LG/L^+G.
\]
\end{definition}

\srcpage{117}

我们将定义 $\mathrm{Sat}_G=\text{「}\Perv(\mathrm{Hk}_G,\Lambda)\text{」}$，并赋予它一个卷积幺半结构。这些都是 $\Rep_{W_E}(\Lambda)$-线性范畴，因为一切都生活在 $\Div^1$ 上。因此 $\Rep_{W_E}(\Lambda)$-线性结构来自与 $\Div^1$ 上局部系统的拉回作张量。

\begin{theorem}
幺半范畴 $(\mathrm{Sat}_G(\Lambda),\star)$ 自然地提升为对称幺半范畴，使得
\[
\oplus H^i(\Gr_{G,\Spa\ol\eta},-)\colon\mathrm{Sat}_G(\Lambda)\to\Rep_{W_E}(\Lambda)
\]
是纤维函子，对应的 Tannaka 群为 $\widehat G$。
\end{theorem}

\begin{remark}
由于 Tannaka 群是 $\Rep_{W_E}$ 中的内蕴对象，它可以看作一个带有自然 $W_E$-作用的约化群 $\widehat G$。
\end{remark}

固定 $\cO_E$ 上的约化群 $G$。还有 Witt 向量仿射 Grassmann 簇 $\Gr_G^{\mathrm{Witt}}$，它是把完美的 $\ol{\Fq}$-代数 $A$ 送到 $\{W_{\cO_E}(A)$ 上、在 $W_{\cO_E}(A)[1/\pi]$ 上平凡化的 $G$-主丛$\}$ 的函子。它带有 $L^+_{\mathrm{Witt}}G$ 的作用，后者是把 $A$ 送到 $G(W_{\cO_E}(A))$ 的函子。我们推出以下结果，它此前已由 Zhu 得到。

\begin{corollary}[Zhu]
我们有对称幺半范畴的等价
\[
(\Perv_{L^+_{\mathrm{Witt}}G}(\Gr_G^{\mathrm{Witt}}),\star)\cong\Rep(\widehat G,\otimes).
\]
\end{corollary}

为推出这一点，我们使用一个退化
\[
\begin{tikzcd}
(\Gr_G^{\mathrm{Witt}})^\diamond \arrow[r] \arrow[d] & \Gr_G^{\mathrm{BD}} \arrow[d] & \Gr_G\times_{\Div^1}\Spa\breveE \arrow[l] \arrow[d] \\
\Spa\ol{\Fq} \arrow[r] & \Spa\cO_{\breveE} & \Spa\breveE \arrow[l]
\end{tikzcd}
\]
于是有一种用 ULA 形式化来「特殊化」反常层的方法。

这个新证明与 Zhu 的证明非常不同。问题在于，对称幺半结构在几何上不容易看出来。在经典情形中，它来自「融合」（fusion）：整体曲线上的两个点相撞。在 $p$ 进情形中，这需要一个像「$\Spec\mathbb Q_p\times\Spec\mathbb Q_p$」那样的空间。在概形的世界里无法以有趣的方式做到这一点，但在菱形的世界里可以做到。
